\chapter{Introduction}
\label{ch:introduction}

\section{Problem}
\label{sec:problem}

A retailer has to decide every day how much of each product to keep in each store. The decision rests on a
forecast of demand: if the forecast is too low, the shelf is empty and the sale is lost; if it is too high,
money is tied up in stock, and perishable goods are thrown away. The forecast is needed at the lowest level,
one product in one store, because that is where replenishment takes place. Forecasts for a whole store, a
product category or a region are needed as well, for purchasing, staffing and planning.

At the lowest level, daily demand looks very different from the smooth, seasonal series of textbooks. Most
products are not sold every day. Their sales are a sequence of zeros interrupted by small positive
quantities, a pattern known as \emph{intermittent demand}. The data used in this thesis, the daily sales of
3{,}049 products in ten Walmart stores published for the M5 forecasting competition
\citep{makridakis2022background}, show how common this is: the typical series has no sale on about 60\% of
the days since the product was introduced, and by the classification of \citet{syntetos2005categorization}
about 91\% of the 30{,}490 series are intermittent or lumpy.

Intermittent demand is hard to forecast for two reasons. The first is that one series carries very little
information. Whether a product is sold on Tuesday or on Thursday is largely a matter of chance, so the
recent history of a single series says little about its next days. The second is that the usual assumptions
of time series models do not hold: the values are small counts with many zeros, not continuous and roughly
normal quantities. Specialised methods were developed for this setting. \citet{croston1972} proposed to
smooth the sizes of the demands and the intervals between them separately; later methods corrected the bias
of this estimator \citep{syntetos2005accuracy}, let the forecast decay during long periods without sales
\citep{teunter2011} or smoothed the demand after aggregating it in time \citep{nikolopoulos2011}. All of
them forecast one series from its own history and produce a flat forecast, the expected demand per day.
They are simple, fast and difficult to beat on single series, but they cannot use what is known in advance
about the days to be forecast (the day of the week, a holiday, a price change), and they learn nothing from
the thousands of other series sold in the same stores.

A \emph{global} model takes the opposite approach: one model is trained on all series together
\citep{januschowski2020,monteromanso2021}. What cannot be learnt from one sparse series, such as the weekly
rhythm of a store or the effect of a promotion, can be learnt from thousands of them. Gradient boosted
decision trees are the most successful family of such models for tabular data. In the M5 competition
most of the leading solutions were global gradient boosting models, trained with a loss function suited
to count data and with a rich set of input features \citep{makridakis2022accuracy}.

That success leaves an open question, because the winning solutions changed several things at once: the
model, the pooling of series, the loss function and the features. It is therefore not clear which of these
choices is responsible for the gain, nor whether the gain exists for the sparse series themselves or only
for the aggregated totals, on which the competition measure puts most of its weight. \citet{tsantilis2025}
separated one of these choices. On 100 intermittent M5 series, and with past sales as the only input, they
compared models trained for each series with models trained on all of them, for three architectures. Pooling
helped the gradient boosting model, but the differences between all the models were small, no classical
method was included for reference, and the authors left untested their hypothesis about what happens when
other inputs are added.

This thesis takes up that question on the complete data set. It studies one model family, gradient boosted
trees as implemented in XGBoost \citep{chen2016}, and varies one design choice at a time: the inputs, the
loss function, the level at which the series are pooled and the strategy for forecasting several days
ahead. Every variant is compared with the classical methods for intermittent demand, both on single series
and on the aggregated levels of the hierarchy.

\section{Research Questions and Contributions}
\label{sec:questions}

The thesis answers the following questions.

\begin{description}[leftmargin=1.5cm,style=nextline]
\item[RQ1] Is a global gradient boosting model more accurate than the classical methods for intermittent
  demand, and is the answer the same for single series as for the aggregated levels of the hierarchy?
\item[RQ2] How much does pooling matter? How do one model per series, one model per store or per demand
  class, and one model for all series compare?
\item[RQ3] Which inputs carry the gain: past sales, inputs known in advance (calendar, events, SNAP days,
  prices), statistics of the sales history, or the identity of the series?
\item[RQ4] Does a loss function designed for count data (Poisson, Tweedie) improve on the squared error?
\item[RQ5] Is it better to forecast each day ahead directly or to apply a one-day-ahead model recursively?
\item[RQ6] Where do the classical methods still win, and what does the model actually use to form its
  forecasts?
\end{description}

The contributions of the thesis are the following.

\begin{itemize}
\item A controlled comparison on all 30{,}490 series of the M5 data, in which the inputs, the loss, the
  level of pooling and the forecasting strategy are varied one at a time, while the data, the tuning
  procedure and the evaluation stay the same.
\item An evaluation with two measures that answer different questions: the mean error of the single series,
  most of which are intermittent, and a revenue-weighted error over seven levels of the hierarchy. The two
  measures give different pictures, and the difference is one of the main findings.
\item A comparison with strong classical methods, including versions whose smoothing constants are chosen
  for each series, so that the gain of the global model is not measured against weak references only.
\item An empirical test of the hypothesis of \citet{tsantilis2025} that additional inputs reduce the
  difference between models per series and a global model.
\item A check of the stability of the conclusions: the main models are retrained with several random seeds
  and evaluated on three earlier periods in addition to the test period, which shows which differences are
  reliable and which change sign from one period to another.
\item An interpretation of the final model with Shapley values, which relates what the model uses to what
  the classical methods use and explains where its advantage comes from.
\end{itemize}

\section{Outline}
\label{sec:outline}

Chapters~\ref{ch:timeseries} to~\ref{ch:evaluation} give the background. Chapter~\ref{ch:timeseries}
introduces time series forecasting, intermittent demand and the classical methods used as references.
Chapter~\ref{ch:boosting} describes gradient boosting, XGBoost, loss functions for data with many zeros and
Shapley values. Chapter~\ref{ch:global} explains how forecasting becomes a supervised learning problem, the
difference between local and global models, the input features and the strategies for forecasting several
steps ahead. Chapter~\ref{ch:evaluation} presents the error measures and the statistical tests used to
compare forecasts.

Chapter~\ref{ch:setup} describes the data, the exploratory analysis and the experimental setup: the split
into training, validation and test periods, the inputs, the reference methods and the tuning of the models.
Chapter~\ref{ch:results} reports the results, one section for each research question, and discusses where
the classical methods remain competitive. Chapter~\ref{ch:conclusion} answers the research questions, states
the limitations of the work and suggests directions for further research.
